\pdfmapfile{=cm.map}
\pdfmapfile{+cmextra.map}
\pdfmapfile{+latxfont.map}
\pdfmapfile{+symbols.map}
\pdfmapfile{+euler.map}
\documentclass[11pt]{article}
\usepackage[margin=1in]{geometry}
\usepackage{amsmath,amssymb,amsthm,booktabs,hyperref,microtype}
\hypersetup{colorlinks=true,urlcolor=blue,linkcolor=blue}
\newtheorem{theorem}{Theorem}[section]
\newtheorem{lemma}[theorem]{Lemma}
\newtheorem{proposition}[theorem]{Proposition}
\newtheorem{corollary}[theorem]{Corollary}
\newcommand{\Q}{\mathbb Q}
\newcommand{\Z}{\mathbb Z}
\newcommand{\F}{\mathbb F}
\newcommand{\HH}{\mathcal H}
\newcommand{\htw}{\operatorname h}
\newcommand{\Norm}{\operatorname N}
\title{Odd elliptic indices and a smaller quartic definition}
\author{Research and exact computations produced with ChatGPT\\for Gabriel Conant}
\date{November 7, 2026}
\begin{document}
\maketitle

\noindent The principal change is an exact height identity on an odd elliptic orbit. It replaces nine height detectors by three and reduces the previous 833-variable construction to 554 variables. A separate formal-group refinement gives 523 variables. All polynomials have degree four, and all counts include the input. These are new arithmetic constructions relative to the inputs of the repository; they are not merely eliminations from its old witness space.

\section{Statement and dependencies}
The reference construction is the direct-target variant of
\href{https://github.com/elliotglazer/diophantine-z-q}{\texttt{elliotglazer/diophantine-z-q}},
at contribution-branch commit
\texttt{772c4adf4a05e21255e4c1784e08e917fda4c68a}.
Its source and arithmetic argument are recorded in
\nolinkurl{research/overnight-20261106/DIRECT_TARGETS.md}.
The main-branch snapshot is
\texttt{5856286a57fba19574a40704571ee70e79738281}.
We retain the point chart and algebraic reductions from the accompanying November~6 note, \emph{A smaller quartic for the H10 construction}.

\begin{theorem}\label{thm:main}
Under the arithmetic inputs of the pinned direct-target construction, the supplied polynomial $P_{554}\in\Z[X,Y_1,\ldots,Y_{553}]$ has total degree four and satisfies
\[
 q\in\Z\quad\Longleftrightarrow\quad
 \exists\mathbf y\in\Q^{553}\ P_{554}(q,\mathbf y)=0
 \qquad(q\in\Q).
\]
Using the additional argument in Section~\ref{sec:depth}, the supplied polynomial $P_{523}\in\Z[X,Y_1,\ldots,Y_{522}]$ has the same property. Both use the fixed coefficient $S_0=3^{2^{26}}$. Replacing this coefficient by a quadratic constant chain gives polynomials with 571 and 540 total variables, respectively, whose coefficients have at most 3,247 bits.
\end{theorem}

The inherited inputs include the generator statement for $2E(K)$, the elliptic height lower bound, the exponent-19 five-contact height estimate, the scalar and anchor lemmas, and the simultaneous-prime and pole-completeness construction. We prove the additional arithmetic assertions below. We do not re-audit the entire repository's proof of these inputs. The 523-variable refinement uses a consequence of the quantitative prime-pattern argument proved in Section~\ref{sec:depth}; it does not assume a stronger prime-pattern theorem.

\begin{center}
\begin{tabular}{lrrr}
\toprule
Construction & Variables & Quadratic residuals & Quartic monomials\\
\midrule
Odd-product construction & 554 & 523 & 33,440\\
With the constant chain & 571 & 540 & 33,490\\
With asymmetric formal depth & 523 & 492 & 35,163\\
With depth and constant chain & 540 & 509 & 35,213\\
\bottomrule
\end{tabular}
\end{center}

The first and third rows are integer polynomials: $S_0$ is a fixed integer coefficient, not a coordinate. Their coefficient tables retain an exact expression in $S_0$. The other rows give every coefficient as a decimal integer. All four tables have primitive content one. The change from 833 to 554 removes 279 variables; the additional depth argument removes 31 more.

\section{An exact product-height identity}\label{sec:height}
Put
\[
 K=\Q(\sqrt2),\quad d=3-\sqrt2,\quad D=d^2=11-6\sqrt2,
 \qquad E:y^2=x^3-Dx,
\]
and put
\[
 R_0=(1+\sqrt2,2\sqrt2),\qquad
 P=2R_0=\left(\frac{57-28\sqrt2}{8},\frac{-660+417\sqrt2}{32}\right).
\]
The source proves that $P$ has infinite order and $2E(K)=\Z P$. Write $\htw$ for the absolute logarithmic Weil height. The following lemma uses the displayed coordinates and doubling identity, together with the infinite order of $P$.

\begin{lemma}\label{lem:product}
For odd integers $n_1,\ldots,n_r$,
\[
 \htw\!\left(\prod_{i=1}^r x([n_i]P)\right)
 =\sum_{i=1}^r\htw(x([n_i]P)).
\]
\end{lemma}
\begin{proof}
Let $T_0=(0,0)$. At a finite place $v$ with $v(D)=0$, use the finite quotient $E(K_v)/E_1(K_v)$, where $E_1$ is the formal subgroup. For a finite nonidentity point $Q$, membership in $E_1$ is equivalent to $v(x(Q))<0$. Since
\[
 x(Q+T_0)=-D/x(Q),
\]
positive valuation of $x(Q)$ is equivalent to $Q+T_0\in E_1$. The image of $T_0$ in this quotient has order two.

If an odd multiple of $P$ belongs to $E_1$, the image of $P$ has odd order. None of its multiples can then equal the image of $T_0$, so all $x([n_i]P)$ have nonpositive valuation. If no odd multiple belongs to $E_1$, all these abscissae have nonnegative valuation. Thus the valuations have a common weak sign. This argument includes the place above two.

The only place with $v(D)>0$ is $\mathfrak d=(3-\sqrt2)$ above seven. The displayed integral model is minimal there, since $v_{\mathfrak d}(-D)=2<4$. Its reduction is $y^2=x^3$ over $\F_7$. Substituting $\sqrt2=3$ shows that $P$ reduces to $(1,6)$, a nonsingular point. All multiples therefore belong to the subgroup $E_0(K_{\mathfrak d})$ of points with nonsingular reduction. A positive valuation of $x$ would force a positive valuation of $y$, and hence reduction to the singular point $(0,0)$. It cannot occur. The valuations again have a common weak sign.

At either real embedding, each $[n_i]P$ is a double of a non-torsion point. The duplication formula makes its abscissa a nonzero square. Its ordinate is nonzero, so the equation of $E$ gives $x([n_i]P)>d$ at that embedding. Both real values $3\pm\sqrt2$ of $d$ exceed one.

At a finite place, numbers whose valuations have the same weak sign satisfy
\[
 \log\max\left(1,\left|\prod_i x([n_i]P)\right|_v\right)
 =\sum_i\log\max(1,|x([n_i]P)|_v).
\]
At both real places the same equality holds because all the factors exceed one. Summing with normalized local weights proves the assertion.
\end{proof}

The product here is an element of $K$. It is not its norm down to $\Q$. Norms can cancel a pole at one split place against a zero at the other, whereas the proof above works separately at every place.

\begin{corollary}\label{cor:scales}
Suppose the four root points $Q_i=[n_i]P$ have odd indices, and write
\[
 C=\prod_{i=1}^4x(Q_i)=C_0+C_1\sqrt2.
\]
The original nine height-detector scales may be replaced by
\[
 1,\qquad C_0^2,\qquad (1+C_1^2)^2.
\]
The root-height estimate retains the original constant $7\HH$, where $\HH$ is the exponent-19 height constant of the source.
\end{corollary}
\begin{proof}
Both real values of $C$ are positive, so $C_0>0$. All three scales are positive nonzero rational squares. If a nonzero integer $N$ is divisible by the $k$th powers of the three contact radicals, the unchanged detector estimate gives
\[
 \htw(\lambda),\ \htw(C_0^2\lambda),\
 \htw((1+C_1^2)^2\lambda)
 \le T:=(\HH+\log3)|N|^{19/k}.
\]
The elementary rational height identities yield
\[
 \htw(C_0)\le T,\qquad \htw(C_1)\le T/2.
\]
Indeed, height subadditivity first bounds the heights of the two scales by $2T$; also $\htw(C_0^2)=2\htw(C_0)$ and $\htw((1+C_1^2)^2)\ge4\htw(C_1)$. Hence
\[
 \htw(C)\le\htw(C_0)+\htw(C_1)+\tfrac32\log2
 \le\tfrac32T+\tfrac32\log2
 <7\HH |N|^{19/k}.
\]
Here $\HH\ge1$. Lemma~\ref{lem:product} bounds each root height by $\htw(C)$. In particular, the inherited lower bound, with $B=\max_i|n_i|$, gives
\begin{equation}\label{eq:height-transfer}
 B^2-6<7\HH |N|^{19/k}.
\end{equation}
The old completeness proof uses its scales only after the roots have been fixed, and requires fixed positive rational squares. Thus its prime selection and pole construction apply to this shorter list.
\end{proof}

\section{Enforcing odd indices by quadratic equations}\label{sec:chart}
Put $\epsilon=1+\sqrt2$. The existing chart for a point $Q$ has $K$-coordinates
\[
 x=h^2,\qquad j^2=x^2-D,\qquad Q=(x,hj).
\]
To this chart add $t,r\in K$ and the equations
\begin{equation}\label{eq:odd-chart}
 t^2=2(x+j),\qquad th+x+j=\epsilon r^2.
\end{equation}
These are four rational quadratic equations in four additional rational coordinates. No extra inverse witnesses are needed.

\begin{proposition}\label{prop:odd-chart}
The projected points of the augmented chart are exactly $P+4E(K)$, that is, the odd multiples of $P$.
\end{proposition}
\begin{proof}
Use the squareclass homomorphism
\[
 \alpha:E(K)\longrightarrow K^*/K^{*2},\qquad
 \alpha(O)=1,\quad \alpha(T_0)=[-D]=[-1],\quad
 \alpha((a,b))=[a]\quad(a\ne0).
\]
For clarity, its homomorphism property follows from the line-intersection identity
\[
 x(R+S)x(R)x(S)=
 \left(\frac{x(R)y(S)-x(S)y(R)}{x(R)-x(S)}\right)^2
\]
when the displayed abscissae are nonzero and distinct. Doubling follows from
$x(2R)=((x(R)^2+D)/(2y(R)))^2$; addition by $T_0$ follows from $x(R+T_0)=-D/x(R)$. Inverse pairs and identity summands are immediate. If $R+S=T_0$, then $x(R)x(S)=-D$. These cases cover the exceptional pairs.

Since $2E(K)=\Z P$ and $2R_0=P$, every point of $E(K)$ has a unique expression $mR_0+T$, where $m\in\Z$ and $T\in E[2](K)$. The two-torsion is generated by $T_0$ and $T_d=(d,0)$, so
\[
 \alpha(mR_0+aT_0+bT_d)=[\epsilon]^m[-1]^a[d]^b.
\]
The classes $[\epsilon],[-1],[d]$ are independent. A square $\epsilon^c(-1)^a d^b$, with $a,b,c\in\{0,1\}$, would have rational square norm $(-1)^c7^b$. This forces $c=b=0$, and the real embeddings then force $a=0$. Consequently
\begin{equation}\label{eq:squareclass}
 \alpha(R)=[\epsilon]\quad\Longleftrightarrow\quad R\in R_0+2E(K).
\end{equation}

Given the augmented chart, set $u=t^2=2(x+j)$ and
\[
 a=th+u/2,\qquad b=at.
\]
Expansion gives
\[
 u(x-u/4)=D,\qquad a^2-ua=D,\qquad b^2=a^3-Da.
\]
In particular $t,a,b$ are nonzero. The tangent slope at $R=(a,b)$ is $h+t$; duplication gives $2R=(x,hj)=Q$. Equation~\eqref{eq:odd-chart} says $a=\epsilon r^2$, so $\alpha(R)=[\epsilon]$. By~\eqref{eq:squareclass}, $Q\in P+4E(K)$.

Conversely, for $Q=(2m+1)P$, take $R=(2m+1)R_0=(a,b)$. Then $a,b\ne0$, and $a=\epsilon r^2$ for some $r\in K$. Set
\[
 t=b/a,\qquad h=(a^2+D)/(2b),\qquad x=h^2,
 \qquad j=t^2/2-x.
\]
The same duplication identities give $Q=(x,hj)$, $j^2=x^2-D$, and both equations~\eqref{eq:odd-chart}.
\end{proof}

Only the four root points receive this augmentation. The shared point and the four comparison points retain their old charts: their indices need not be odd. The total cost is sixteen coordinates and sixteen quadratic residuals. The projected-point formulation matters, since the added equations can select a sign representation of the old $h,j$ chart.

\section{Odd labels, scalar bounds, and completeness}\label{sec:labels}
Write the root points as
\[
 A=[n]P,\qquad S=[a]P,\qquad T=[t]P,\qquad Z=[b]P.
\]
Here $S$ denotes a point; the fixed coefficient remains $S_0$. Replace the old input expression by
\[
 \beta=2X^2+3.
\]
For rational $X$, integrality of $\beta$ is equivalent to integrality of $X$: if $X=u/v$ in lowest terms, then $\beta\in\Z$ forces $v^2\mid2$, and hence $v=1$. For integral $X$, $\beta$ is odd.

In the scalar norm block and the physical anchor, replace $S$ by
\[
 R=S-P=[a-1]P.
\]
The comparison and product-height blocks still use $S$. The retained subtraction guard excludes $a=\pm1$, which is harmless for completeness. The original scalar contact product is
\[
 \beta(\beta-1)(\beta-2401)(32\beta+343)(2\beta-1).
\]
It never vanishes for rational $X$: $\beta\ge3$, and $\beta=2401$ would give $X^2=1199$, which is not a rational square. At an odd prime, a negative valuation of $\beta$ is twice that of $X$, as required by the scalar argument.

\subsection{Soundness}
Put $c=|a-1|$. Applying the scalar and anchor lemmas to $R$ gives
\[
 \htw(\beta)\le\HH c^{19/32},\qquad
 \log_2c>25{,}165{,}823.5,
 \qquad 1\le\HH<2^{2{,}093{,}966}.
\]
The anchor is still $S_0=3^{2^{26}}$. These explicit bounds imply
\[
 c>(80\HH)^{10},\qquad
 |a|\ge c-1\ge c/2>(40\HH)^{10},
\]
and
\begin{equation}\label{eq:margins}
 \htw(\beta)<c/80\le |a|/40,
 \qquad\log(40\HH)<c/80\le |a|/40.
\end{equation}
For example $10(7+2{,}093{,}966)<25{,}165{,}823.5$, and $80<2^7$. The first estimate in~\eqref{eq:margins} follows from $c^{13/32}>80\HH$; the logarithmic estimate is weaker still.

The local comparison lemmas retain their original conclusions. Let
\[
 F=t^2-a(a-2)tb+(a-2)^2b^2-1,
 \qquad B=\max(|n|,|a|,|t|,|b|).
\]
For $B>4$, the triangle inequality gives $|F|\le2B^4+6B^3+5B^2+1<5B^4$. If $F\ne0$, the three radicals divide it to order 40, so~\eqref{eq:height-transfer} gives
\[
 B^{1/10}<7\HH\,5^{19/40}+6<24\HH.
\]
This contradicts the anchor. Thus $F=0$. Since the point $Z$ is finite and nonidentity, $b\ne0$. The source's Lucas growth lemma now gives $\log B>|a|-3$.

If $n\ne\beta$, the nonzero numerator $N$ of $n-\beta$, in reduced form, satisfies $|N|\le2e^{\htw(\beta)}B$, and the radicals divide it to order 12. Equation~\eqref{eq:height-transfer} and $B>4$ give
\[
 B^{5/12}<40\HH e^{(19/12)\htw(\beta)},
\]
so
\[
 \log B<\frac{12}{5}\log(40\HH)+\frac{19}{5}\htw(\beta)
 <\frac{31}{200}|a|<|a|-3.
\]
This is the other contradiction. Therefore $n=\beta\in\Z$, and $X\in\Z$.

\subsection{Completeness}
Fix $X\in\Z$, and set $n=\beta$. Perform the old scalar preparation: choose its finite correction and then a positive shared index $h_0$ divisible by 64 and sufficiently deep at the fixed small primes. Choose a positive scalar index $s$ sufficiently divisible to satisfy the scalar leaves and physical anchor for $[s]P$. We may require $6\mid s$ and take $s$ arbitrarily large. Set $a=s+1$. Then $a$ is odd, $a\equiv1\pmod3$, and $S-P=[s]P$ is exactly the prepared scalar point.

Define
\[
 U_0=0,\qquad U_1=1,\qquad U_{j+1}=aU_j-U_{j-1},
 \qquad k=2(a-2),
\]
and take $b=U_k/(a-2)$ and $t=U_{k+1}$. Modulo $a-2$, induction gives $U_j\equiv j$, so $b\in\Z$. Modulo two the sequence has period $0,1,1$; since $k\equiv1\pmod3$, both $U_k$ and $U_{k+1}$ are odd. The positive odd denominator $a-2$ shows that $b$ is odd. The Lucas identity gives
\[
 t^2-a(a-2)tb+(a-2)^2b^2=1.
\]
Thus all four root indices $n,a,t,b$ are positive and odd. Proposition~\ref{prop:odd-chart} supplies their chart witnesses.

Fix these roots and the three scales of Corollary~\ref{cor:scales}. Local continuity of the scalar and comparison functions at the prepared shared point gives the same fixed modulus $L$, with $64\mid L$. Only then choose the contact primes by the inherited simultaneous-prime theorem. The direct-target construction chooses a common index $h$ satisfying the primary and opposite reduction-order congruences and $h\equiv h_0\pmod L$, with the required real bounds. Its comparison points are $H=[h]P$ and $U_i=[hn_i]P$. Their shifted points have indices $h-1$ and $(h-1)n_i$.

All quotient-domain arguments remain valid. The product scales are fixed positive squares, so the original pole witnesses apply unchanged. The scalar block is preserved by the congruence modulo $L$. This gives a rational solution for every integral $X$ and completes the arithmetic proof of the primary construction.

\section{Replacing five logarithms by asymmetric formal depth}\label{sec:depth}
This section is used only for the 523-variable construction. Let
\[
 U_0=H=[h]P,\quad U_i=[u_i]P,\quad
 V_0=H-P,\quad V_i=U_i-Q_i,
\]
where $Q_i=[n_i]P$ are the four roots. Put $f=1/x$ and $z=-x/y$. Keep the alternate truncated logarithms $\ell_{40}(V_i)=z(V_i)G_9(Df(V_i)^2)$, but replace every primary truncated logarithm by $z(U_i)$. Thus
\[
 L_i=z(U_i)/z(H),\qquad
 L_i'=\overline{\ell_{40}(V_i)/\ell_{40}(V_0)}.
\]
The bar denotes conjugation in $K$. The sixteen tests now have the following orders:
\[
\begin{array}{c|l}
1 & f(U_1),\ldots,f(U_4),\ \overline{f(V_0)},\ldots,\overline{f(V_4)}\\[2pt]
40 & f(H),\ L_1-L_1',\ldots,L_4-L_4',\ L_1-\beta,\
       \text{the direct Lucas residual}.
\end{array}
\]
The scalar tests retain their original order 32. The stronger beta tests imply all former order-12 conclusions.

\subsection{Local soundness}
At a consumed unramified place above $p>76176$, a formal point $Q$ with $r=v(z(Q))\ge1$ satisfies
\begin{equation}\label{eq:zerror}
 \log_E(Q)/z(Q)\in1+p^{4r}\mathcal O.
\end{equation}
This is the same formal-logarithm expansion used in the source, starting its tail at the first nonconstant term. Namely $v(f(Q))=2r$, and for its coefficients $a_j$ the factorial estimate is $v(a_j)\ge-(4j+1)/(p-1)$. If $4j+1<p$, the coefficient is integral. For the remaining $j\ge19044$, the displayed bound still makes the $j$th term have valuation at least $4r$. Both $z(Q)$ and its full logarithm have valuation $r$.

The new test on $f(H)$ gives $v(z(H))\ge20$. The order-one tests put the other points on the formal branch. The full logarithm ratios at the primary and conjugate places are
\[
 r_i=u_i/h,\qquad r_i'=(u_i-n_i)/(h-1).
\]
At least one of $h,h-1$ is a $p$-unit, so at least one of these ratios is integral. Even before knowing the depth of $U_i$, the weak unit estimate~\eqref{eq:zerror} and the alternate truncation estimate show that $L_i,L_i'$ have the same valuations as $r_i,r_i'$. The order-40 cross test transfers integrality to both sides. Now $L_i$ integral implies $v(z(U_i))\ge v(z(H))\ge20$. The primary logarithmic corrections are therefore $1$ modulo $p^{80}$. The alternate corrections are $1$ modulo $p^{40}$. Since the ratios are integral, these relative estimates become absolute:
\[
 L_i\equiv r_i\equiv r_i'\pmod {p^{40}}.
\]
Multiplying by $h$ and $h-1$ and subtracting gives $L_i\equiv n_i\pmod {p^{40}}$. At an inert prime conjugation preserves the local valuation, so the same proof applies. The height contradictions of Section~\ref{sec:labels} remain unchanged.

\subsection{A stronger choice within the existing prime-pattern argument}
For completeness, the common index must acquire extra primary divisibility. The following observation ensures compatibility with the opposite reduction orders.

\begin{lemma}\label{lem:separation}
The contact primes in the source completeness argument may be chosen so that no selected prime divides any opposite reduced group order.
\end{lemma}
\begin{proof}
After merging repeated contact forms and fixing all scale and congruence data, the prime-pattern proof uses distinct forms $p_j=u+c_jw^2$, with fixed coefficients $c_j$. Integer dilations and ray congruences change only fixed constants. It counts $u\in[N/2,N]$ and $1\le w\le T$, where $T=\lfloor\tau(N)\sqrt N\rfloor$ and $\tau(N)\to0$ can be chosen arbitrarily slowly in the final diagonal argument. Its weighted count, after the prescribed obstruction removal, retains a fixed positive proportion of $NT$. These are the fixed-field weighted count and final diagonal argument in \nolinkurl{manuscripts/quartic_v8.tex} at the pinned commit.

If there are $k$ forms, impose also
\[
 \frac{\tau(N)N^{1/6}}{(\log N)^k}\longrightarrow\infty.
\]
This is compatible with the same diagonal choice: for each fixed positive value of $\tau$, its threshold can be added to the existing finite list of thresholds. The strip $w<N^{1/3}$ has at most $O(N^{4/3})$ integer pairs and weight at most $O((\log N)^k)$ per pair. Its total weight is $o(NT)$, so an admissible pattern remains outside it.

For these patterns $w^2\ge N^{2/3}$ and $w^2=o(N)$. All $p_j$ are asymptotic to the same $u$. Hasse's bound gives $d_+(p_j)=p_j+O(\sqrt N)$ for the opposite reduced group order. Every ratio $d_+(p_j)/p_i$ is eventually strictly between $1/2$ and $2$. Hence $p_i\mid d_+(p_j)$ would force equality. If $i\ne j$, this contradicts
\[
 |p_i-p_j|\ge c_*N^{2/3},\qquad
 c_*:=\min_{i\ne j}|c_i-c_j|>0.
\]
If $i=j$, equality is impossible because $4\mid d_+(p_i)$ whereas $p_i$ is odd.
\end{proof}

\subsection{Completeness with deeper primary points}
Keep the choice order in Section~\ref{sec:labels}: the roots, scales and continuity modulus $L$ precede the prime selection. Choose primes using Lemma~\ref{lem:separation}, beyond all prime factors of $L$ and the fixed scale data. Put
\[
 M_\pm=\mathop{\mathrm{lcm}}_p(d_\pm(p)/4),\qquad \Pi=\prod_p p.
\]
The direct-target prime theorem gives
\[
 \gcd(M_-,M_+)=1,\quad
 \gcd(M_-,L)=16,\quad\gcd(M_+,L)=1.
\]
Lemma~\ref{lem:separation} gives $\gcd(\Pi,M_+)=1$, and the cutoff gives $\gcd(\Pi,L)=1$. Thus
\[
 h\equiv0\pmod{\Pi^{19}M_-},\qquad
 h\equiv1\pmod{M_+},\qquad h\equiv h_0\pmod L
\]
are compatible, since $64\mid h_0$. The inherited real-progression argument supplies a sufficiently large $h$ in this progression with the same real bounds.

The point $P$ is globally a double, so multiplication by the primary doubling-image order first makes it formal, with $z$-valuation at least one. Multiplication by $p^{19}$ raises this valuation by nineteen, using the formal logarithm isometry. Consequently $H=[h]P$ and $U_i=[hn_i]P$ have primary depth at least twenty, and their parameter ratios equal $n_i$ to order eighty. At the opposite places, $V_0=[h-1]P$ and $V_i=[(h-1)n_i]P$ are formal and their retained logarithm ratios equal $n_i$ to order forty.

The pole-completeness estimates still apply to these fixed test orders. In the source notation their large-real-parameter decay is $\lambda^{-3}$ for each order $m$, and the small-prime lower bound is $3\nu+(10m-8)\eta-2c$, which still tends to infinity with the preparatory divisibility. Scalar conditions and the anchor are preserved by $h\equiv h_0\pmod L$. This proves completeness for the optional refinement.

\section{Polynomial construction and exact verification}
The generators apply three changes to the earlier point-chart source: the four odd-coset charts, the product and three scales, and the shifted scalar point with $\beta=2X^2+3$. The depth variant first changes the logarithms and test orders as in Section~\ref{sec:depth}. Each field equation is represented by two rational component equations. Intermediate products are assigned coordinates as needed to keep every residual quadratic.

The full-log source has 577 raw coordinates; the pinned upstream compiler reduces this to 561. Seven further certified eliminations give 554. The depth source has 540 raw coordinates; compilation gives 532 and nine certified eliminations give 523. Every elimination is justified by exact polynomial identities and reversible substitutions preserving the input. Clearing residual denominators and taking the sum of their squares gives the displayed integer quartics. Over $\Q$, the sum vanishes exactly when every residual does.

The constant-chain versions add $c_0,\ldots,c_{16}$ with
\[
 c_0=3^{1024},\qquad c_{i+1}=c_i^2\quad(0\le i<16).
\]
These witnesses have unique rational values and $c_{16}=S_0$. Substitution for $S_0$ leaves every residual quadratic. Direct expansion gives the 3,247-bit maximum coefficient bound.

The accompanying archive contains the complete polynomial tables, residual systems, coordinate maps, elimination certificates, pinned source generators, arithmetic notes and independent checkers. Each final package has a standard-library-only \texttt{verify\_all.py}. It replays all eliminations, verifies the unique constant lift, independently expands the sums of squares, and compares every coefficient. Both full packages have passed fresh independent runs. A separate source replay rebuilds the two compiled precursors and checks the match to their frozen residual systems.

These finite checks certify the counts, degrees and stated algebraic transformations. The arithmetic assertions in Sections~\ref{sec:height}--\ref{sec:depth} received separate adversarial AI review; they have not received independent human peer review or a complete formal proof check. The integer-projection conclusion remains relative to the pinned repository's arithmetic inputs.

\end{document}
